% Appendix A -- Derivations

\chapter{Derivations}

\label{AppendixA}

This appendix collects the derivations referenced in Chapter~\ref{chap3:methodology}.

\section{Richardson extrapolation coefficients and variance inflation}
\label{sec:appA_richardson}

Let $\mu(\lambda)$ be the expectation value measured at noise scale $\lambda$ and
suppose it is smooth in $\lambda$, so that near zero it admits the expansion
\begin{equation}
    \mu(\lambda) \;=\; \mu_{\text{ideal}} \;+\; \sum_{j=1}^{\infty} a_j \lambda^{j}.
    \label{eq:appA_expansion}
\end{equation}
Given measurements at $m+1$ distinct scales $\lambda_0, \ldots, \lambda_m$, the
unique degree-$m$ interpolating polynomial evaluated at $\lambda = 0$ is, by the
Lagrange form,
\begin{equation}
    \widehat{\mu}_{\text{ZNE}} \;=\; \sum_{k=0}^{m} c_k \, \mu(\lambda_k),
    \qquad
    c_k \;=\; \prod_{j \neq k} \frac{0 - \lambda_j}{\lambda_k - \lambda_j}
    \;=\; \prod_{j \neq k} \frac{\lambda_j}{\lambda_j - \lambda_k}.
    \label{eq:appA_coeffs}
\end{equation}
Substituting eqn.~\eqref{eq:appA_expansion} into eqn.~\eqref{eq:appA_coeffs} and using
the interpolation identities
$\sum_k c_k = 1$ and $\sum_k c_k \lambda_k^{\,j} = 0$ for $1 \leq j \leq m$, the
first $m$ orders of the noise expansion cancel exactly and
\begin{equation}
    \widehat{\mu}_{\text{ZNE}} \;=\; \mu_{\text{ideal}}
    \;+\; a_{m+1} \left( \sum_{k=0}^{m} c_k \lambda_k^{\,m+1} \right)
    \;+\; O(\lambda^{m+2}),
    \label{eq:appA_bias}
\end{equation}
so the residual bias is of order $m+1$ in the noise strength. This is the entire
justification for the method: increasing the polynomial order cancels more terms.

The cost appears when the measurements are noisy. With the three
parameter-scaling coordinates $\lambda \in \{1,2,3\}$ and $m=2$,
eqn.~\eqref{eq:appA_coeffs} gives
\begin{equation}
    c_0 = \frac{(2)(3)}{(2-1)(3-1)} = 3,
    \qquad
    c_1 = \frac{(1)(3)}{(1-2)(3-2)} = -3,
    \qquad
    c_2 = \frac{(1)(2)}{(1-3)(2-3)} = 1,
    \label{eq:appA_values}
\end{equation}
which indeed sum to unity. If each measurement carries independent statistical
variance $\sigma^2$, the estimator variance is
\begin{equation}
    \mathrm{Var}\!\left[ \widehat{\mu}_{\text{ZNE}} \right]
    \;=\; \sigma^2 \sum_{k} c_k^{2}
    \;=\; \sigma^2 \left( 9 + 9 + 1 \right)
    \;=\; 19\,\sigma^2 ,
    \label{eq:appA_variance}
\end{equation}
an inflation of more than an order of magnitude relative to a single measurement.
Richardson extrapolation therefore trades bias for variance, and at a finite shot
budget the trade need not be favourable -- which is the mechanism behind the
shot-budget results of Section~\ref{sec:ch5_shots}. Least-squares fitting of a
degree-one polynomial to the same three points, as used by the linear arm, has
smaller coefficient magnitudes and therefore less variance inflation, at the price of
cancelling only the first order of eqn.~\eqref{eq:appA_expansion}.
For folding, the actual layer ratios determine the weights. At depths
$3$ and $5$ the corresponding equal-variance factors are $21.34$ and
$21.72$, respectively.

\section{Why exponential extrapolation is motivated, and why it fails here}
\label{sec:appA_exp}

Under a global depolarising channel of strength $\varepsilon$ acting on an
$n$-qubit state,
\begin{equation}
    \mathcal{E}(\rho) \;=\; (1 - \varepsilon)\rho \;+\; \varepsilon \frac{I}{2^{n}},
    \label{eq:appA_global}
\end{equation}
any traceless observable $O$ has $\mathrm{Tr}[O \, I/2^n] = 0$, so
$\langle O \rangle$ is simply attenuated by $(1 - \varepsilon)$. If each of $L$
circuit layers applies such a channel independently, the attenuation compounds to
$(1-\varepsilon)^{L} = e^{L \ln (1 - \varepsilon)}$, which is exponential in circuit
depth and motivates fitting
$\mu(\lambda) = a e^{-b \lambda} + c$.

The difficulty is one of parameter counting rather than of physics. With $K = 3$
noise scales and three free parameters $(a, b, c)$ the fit is exactly determined:
there are no residual degrees of freedom, no averaging of the statistical
fluctuations, and the covariance of the fit is not estimable. A fluctuation in any
one of the three measurements propagates directly into the extrapolated value, and
because the model is non-linear in $b$ the propagation can be violent. The results
of Section~\ref{sec:ch5_main} show this arm performing worse than no mitigation at
all, which is a statement about the three-point exactly-determined fit, not about
the exponential ansatz. Four or more noise scales would make the fit
overdetermined and the arm usable.

\section{Fidelity kernel and concentration}
\label{sec:appA_kernel}

For a feature map $x \mapsto \ket{\phi(x)} = U(x)\ket{0}^{\otimes n}$ implemented by
a parametrised circuit, the fidelity kernel is
\begin{equation}
    k(x, x') \;=\; \left| \braket{\phi(x) | \phi(x')} \right|^{2}
    \;=\; \left| \bra{0}^{\otimes n} U^{\dagger}(x) U(x') \ket{0}^{\otimes n} \right|^{2},
    \label{eq:appA_fidelity}
\end{equation}
which is symmetric, satisfies $k(x,x) = 1$, and is positive semi-definite because it
is a Gram matrix of unit vectors under the Hilbert--Schmidt inner product. Hence it
is a valid kernel and may be used in any kernel method. In simulation the states
$\ket{\phi(x_i)}$ can be computed once and stored as rows of a matrix $\Psi$, after
which the entire Gram matrix is the single product
$K = \left| \Psi \Psi^{\dagger} \right|^{2}$ taken elementwise -- the reason the
kernel arm is computationally cheap in this study but would require
$O(M^2)$ circuit evaluations on hardware for $M$ training points.

The concentration phenomenon follows from a dimensional argument. If the embedding
spreads states so that they are close to Haar-random in a $2^n$-dimensional space,
the expected fidelity between two independent states is
\begin{equation}
    \mathbb{E}\left[ \left| \braket{\phi | \psi} \right|^{2} \right]
    \;=\; \frac{1}{2^{n}},
    \label{eq:appA_haar}
\end{equation}
which is exponentially small in the qubit count. The Gram matrix then approaches the
identity: every off-diagonal entry is near zero while the diagonal is exactly one.
A kernel ridge regressor built on such a matrix can interpolate its training data
but has almost no ability to relate a test point to any training point, so its
predictions collapse towards the mean and its test error approaches that of a
constant model -- or worse, since the small residual off-diagonal structure is
dominated by sampling noise. Rescaling the input by a bandwidth
$\beta < 1$ before embedding contracts the embedded states into a smaller region,
restores non-trivial off-diagonal structure, and thereby controls the effective
smoothness of the induced function class. This is why bandwidth is treated in
Section~\ref{sec:ch3_hybrid} as a hyperparameter that must be tuned rather than a
detail that may be left at its default, and why
Section~\ref{sec:ch5_bandwidth} reports the untuned case as a finding.

\section{Why a trace-distance revival implies memory}
\label{sec:appA_blp}

The criterion used in Section~\ref{sec:ch3_blp} rests on one property of quantum
channels. For any completely positive trace-preserving map $\Lambda$ and any pair of
states,
\begin{equation}
    D\!\left( \Lambda(\rho_1), \Lambda(\rho_2) \right)
    \;\leq\; D\!\left( \rho_1, \rho_2 \right),
    \qquad
    D(\rho_1, \rho_2) = \tfrac{1}{2}\left\| \rho_1 - \rho_2 \right\|_1,
    \label{eq:appA_contraction}
\end{equation}
that is, channels are contractive in trace distance. Operationally $D$ is the
optimal bias with which a single measurement can distinguish the two states, so
eqn.~\eqref{eq:appA_contraction} says no channel can make two states easier to tell
apart.

Suppose the layer-to-layer dynamics is Markovian in the CP-divisible sense: the map
from layer $l$ to layer $l+1$ is itself a channel $\Lambda_{l+1,l}$, independent of
history. Applying eqn.~\eqref{eq:appA_contraction} to that map gives
\begin{equation}
    D^{(l+1)} \;=\; D\!\left( \Lambda_{l+1,l}\rho_1^{(l)},\,
    \Lambda_{l+1,l}\rho_2^{(l)} \right) \;\leq\; D^{(l)}
    \label{eq:appA_monotone}
\end{equation}
for every $l$ and every initial pair, so $D^{(l)}$ is non-increasing and the sum of
positive increments in eqn.~\eqref{eq:ch3_blp} vanishes identically. A measured
revival $D^{(l+1)} > D^{(l)}$ is therefore incompatible with CP-divisibility: the
step from $l$ to $l+1$ cannot be a channel acting on the system alone, which is to
say the environment has returned information it previously held. This is why
$\mathcal{N} > 0$ is evidence of memory and not merely of strong noise --- noise of
any strength, provided it is Markovian, can only decrease $D$.

The converse does not hold: $\mathcal{N} = 0$ does not prove CP-divisibility, since
the maximisation over state pairs is taken over a finite set in practice and some
non-Markovian maps have no trace-distance revival at all. The measure is therefore
used here as a positive witness of memory in the production model and as a null
check on the control, not as a proof of Markovianity for the
control~\cite{rivas2014nonmarkov}.

\section{Noise amplification by unitary folding}
\label{sec:appA_folding}

Folding amplifies noise without changing the ideal action of the circuit. Let a
layer's intended unitary be $L$ and let the noise accompanying one execution of it
be a channel $\mathcal{E}$, so the physical layer implements
$\mathcal{L} = \mathcal{E} \circ \mathcal{U}_L$ with
$\mathcal{U}_L(\rho) = L \rho L^{\dagger}$. Replacing $L$ by $L L^{\dagger} L$
executes three physical layers,
\begin{equation}
    \mathcal{L}_{\text{fold}} \;=\;
    \left( \mathcal{E} \circ \mathcal{U}_{L} \right) \circ
    \left( \mathcal{E} \circ \mathcal{U}_{L^{\dagger}} \right) \circ
    \left( \mathcal{E} \circ \mathcal{U}_{L} \right),
    \label{eq:appA_fold}
\end{equation}
whose ideal part is $\mathcal{U}_L \mathcal{U}_{L^\dagger} \mathcal{U}_L =
\mathcal{U}_L$, unchanged, while three noise channels act instead of one. For a
depolarising channel of strength $p$ the composition of $k$ copies has strength
$1 - (1-p)^{k} = kp + O(p^2)$, so to first order in $p$ the effective noise scale of
a folded layer is $3$. Folding $j$ of the $d$ layers, a fraction $f=j/d$, gives the
executed-to-original layer ratio
\begin{equation}
    \lambda_{\text{fold}} \;=\; 1 + 2f,
    \label{eq:appA_lambda}
\end{equation}
which the benchmark uses directly as its extrapolation coordinate. The middle
setting takes $j=\operatorname{round}(d/2)$, yielding a depth-dependent
coordinate for odd $d$ (Section~\ref{sec:ch3_folding}).

The identification of the layer ratio with a noise-strength multiplier
assumes identical depolarising noise on each layer and is accurate to first
order in $p$. Coherent errors depend on the sequence of layer operations.
The defect also carries its state between layers in the production model, and
the always-on crosstalk acts during every extra layer.
The realised layer ratio therefore fixes the extrapolation coordinate
while the response to that coordinate depends on the noise model.

\section{Circuit cost of a shot-estimated fidelity kernel}
\label{sec:appA_kernelcost}

The simulator computes each fidelity entry as the squared magnitude of an inner
product of state vectors. To model hardware acquisition, each entry
\begin{equation}
    k_{ij} \;=\; \left| \left\langle \phi(x_i) \middle| \phi(x_j) \right\rangle \right|^{2}
    \label{eq:appA_kij}
\end{equation}
must be measured, for instance by the inversion test, which prepares
$U^{\dagger}(x_j) U(x_i) \ket{0}$ and estimates the probability of the all-zeros
outcome. For an ideal overlap circuit and independent shots, the $M$-shot estimator is
$\hat{k}_{ij} = B/M$ with $B \sim \mathrm{Binomial}(M, k_{ij})$, so
\begin{equation}
    \mathbb{E}\!\left[\hat{k}_{ij}\right] = k_{ij},
    \qquad
    \mathrm{Var}\!\left[\hat{k}_{ij}\right] = \frac{k_{ij}(1 - k_{ij})}{M}.
    \label{eq:appA_kvar}
\end{equation}
Let $T$ denote the training set at a given fit and $U$ denote the points to
predict. Diagonalising the measured training block gives
$\widehat K_{TT}=V\Lambda V^\top$. Keeping eigenvalues above
$|T|\epsilon_{\mathrm{mach}}\max_i|\lambda_i|$ yields
$K_{TT,+}=V_+\Lambda_+V_+^\top$. The cross-kernel is projected onto the retained training subspace,
$K_{UT,+}=\widehat K_{UT}V_+V_+^\top$. With ridge penalty $\alpha>0$, the
predictions are
\begin{equation}
    \widehat y_U
    = \widehat K_{UT} V_+
      (\Lambda_+ + \alpha I)^{-1} V_+^\top y_T.
    \label{eq:appA_shot_projection}
\end{equation}
The same operation applies to each inner validation set and to the final test
set, with the respective training block defining $V_+$.

Training on $n$ points and predicting on $m$ points uses a symmetric training
block and a rectangular cross-kernel. Counting the diagonal, their cost is
\begin{equation}
    C(n,m) = M \left[ \frac{n(n+1)}{2} + nm \right]
    \label{eq:appA_cost}
\end{equation}
circuit executions. In R4 each grouped fold fits its feature scaler on its own
training rows, so the simulated acquisition cost of tuning is the sum of
$C(n_f,m_f)$ across the four folds. The matrices within a fold are shared across
all tested ridge penalties. The final fit and held-out evaluation add
$C(3150,1350)$, equivalent to $9{,}215{,}325$ entries before multiplying by
$M$. A subsequent prediction adds $3150M$ fidelity shots. Acquisition of the
three input measurements and labelled reference data is accounted for separately.


\section{Finite-sampling model}
\label{sec:appA_shots}

Let $O$ be a Pauli string, so $O^2 = I$ and its eigenvalues are $\pm 1$. Writing
$q$ for the probability of the $+1$ outcome and $\mu$ for the expectation value,
\begin{equation}
    \mu \;=\; (+1)\,q \;+\; (-1)(1-q) \;=\; 2q - 1
    \qquad \Longleftrightarrow \qquad
    q \;=\; \frac{1 + \mu}{2}.
    \label{eq:appA_q}
\end{equation}
An $N$-shot experiment observes $k \sim \mathrm{Binomial}(N, q)$ outcomes equal to
$+1$, and the natural estimator is $\widehat{\mu}_N = 2k/N - 1$. It is unbiased,
$\mathbb{E}[\widehat{\mu}_N] = 2q - 1 = \mu$, with variance
\begin{equation}
    \mathrm{Var}\!\left[ \widehat{\mu}_N \right]
    \;=\; \frac{4}{N^{2}} \, N q (1-q)
    \;=\; \frac{1 - \mu^{2}}{N},
    \label{eq:appA_var}
\end{equation}
so the sampling noise is largest for observables whose expectation is near zero and
vanishes as $|\mu| \to 1$. Applying eqn.~\eqref{eq:appA_q} to the simulator expectation reproduces the
sampling distribution for independent Pauli shots under the simulated measurement
model. Storing both the exact and the sampled values additionally makes the
shot-budget sweep of Section~\ref{sec:ch5_shots} a re-sampling operation rather than
a re-simulation, since eqn.~\eqref{eq:appA_var} can be realised at any $N$ from the
same exact values.
